\documentclass[11pt]{article}
\usepackage[T1]{fontenc}
\usepackage{textcomp}
\usepackage[margin=0.75in]{geometry}
\usepackage{amsmath,amssymb,amsthm}
\usepackage{array,booktabs}
\usepackage{hyperref}
\setlength{\parindent}{0pt}
\newtheorem{theorem}{Theorem}
\theoremstyle{definition}
\newtheorem{definition}{Definition}
\title{Flow Proof Artifact: prop-05-isosceles-base-angles}
\author{Generated by \texttt{flow doc proof}}
\date{Flow Proof Book}
\begin{document}
\maketitle
In an isosceles triangle, the angles at the base equal one another.\\[0.5em]
\textbf{Source.} euclid — Elements, Book I, Proposition 5\\[0.5em]
\bigskip
\noindent{\large \textbf{Derived fact 1.} \textit{In an isosceles triangle, the angles at the base equal one another} \hfill the Euclidean plane \cdot Euclid Book I \cdot Proposition 5: base angles of an isosceles triangle are equal}\par
\smallskip\noindent\textit{Coordinate: the Euclidean plane \cdot Euclid Book I \cdot Proposition 5: base angles of an isosceles triangle are equal}\par
\smallskip\noindent\textit{Source: euclid — Elements, Book I, Proposition 5}\par
\smallskip\noindent\textit{Built on: proposition 4: side-angle-side congruence, for Book I of the Elements in on the Euclidean plane}\par
\medskip
\noindent\textbf{Goal.} In an isosceles triangle, the angles at the base equal one another\par
\noindent$\displaystyle \angle B = \angle C$\par
\medskip
\noindent
\begin{tabular}{@{} >{\raggedright\arraybackslash}p{0.06\textwidth} >{\raggedright\arraybackslash}p{0.47\textwidth} >{\raggedleft\arraybackslash}p{0.41\textwidth} @{}}
\textbf{\#} & \textbf{Proof} & \textbf{Mathematics} \\
\midrule
\textbf{1.} & We prove that proposition 5: base angles of an isosceles triangle are equal for Book I of the Elements in the Euclidean plane. & \\[0.45em]
\textbf{2.} & Let D = point\_on\_BC\_extended\_with\_BD\_equal\_to\_CE. & \\[0.45em]
\textbf{3.} & Let E = point\_on\_CB\_extended\_with\_CE\_equal\_to\_BD. & \\[0.45em]
\textbf{4.} & We invoke the derived fact: segment\_AB == segment\_AC. & \\[0.45em]
\textbf{5.} & We invoke the axiom governing Book I of the Elements in the Euclidean plane: proposition 4: side-angle-side congruence, for Book I of the Elements in on the Euclidean plane. & \\[0.45em]
\textbf{6.} & From step 2, step 3, step 4, and step 5, we can deduce that angle ABD equals angle ACE. & $\displaystyle \angle ABD = \angle ACE$ \\[0.45em]
\textbf{7.} & We can deduce that angle ABC equals angle ACB. & $\displaystyle \angle ABC = \angle ACB$ \\[0.45em]
\textbf{8.} & We can deduce that angle at B equals angle at C. Hence proven. & $\displaystyle \angle B = \angle C$ \\[0.45em]
\bottomrule
\end{tabular}
\label{thm:Geometry-Euclid Book I-proposition_5_base_angles_of_an_isosceles_triangle_are_equal}
\par\medskip\hrule
\end{document}
